\documentclass[a4paper,12pt]{article}
\usepackage[T2A]{fontenc}
\usepackage[utf8]{inputenc}
\usepackage[russian]{babel}
\usepackage{amsmath,amssymb,mathtools}
\usepackage{helvet}
\renewcommand{\familydefault}{\sfdefault}
\usepackage{icomma}
\usepackage[a4paper,top=23mm,bottom=23mm,left=24mm,right=24mm]{geometry}
\usepackage{microtype}
\usepackage{xcolor}
\usepackage{tikz}
\usetikzlibrary{calc,arrows.meta,positioning,shapes.geometric}
\usepackage{pgfplots}
\pgfplotsset{compat=1.18}
\usepackage{booktabs,tabularx}
\usepackage{enumitem}
\usepackage{fancyhdr}
\usepackage{setspace}
\usepackage{hyperref}
\definecolor{accent}{HTML}{008080}
\definecolor{deep}{HTML}{005B5B}
\definecolor{textgray}{HTML}{404648}
\definecolor{second}{HTML}{B66440}
\hypersetup{colorlinks=true,linkcolor=deep,urlcolor=deep,pdfauthor={Лёвин Артём Александрович},pdftitle={Чек-лист: квадратичная функция и графики}} 
\setlength{\parindent}{0pt}
\setlength{\parskip}{6pt}
\onehalfspacing
\setlist[enumerate]{leftmargin=7mm,itemsep=3pt,topsep=3pt}
\pagestyle{fancy}
\fancyhf{}
\fancyhead[L]{\small\color{textgray}Квадратичная функция и графики}
\fancyhead[R]{\small\color{textgray}Ксения Клыкова}
\fancyfoot[C]{\small\color{textgray}\thepage}
\renewcommand{\headrulewidth}{0pt}
\renewcommand{\footrulewidth}{0pt}
\setlength{\headheight}{16pt}
\newcommand{\lessonsec}[1]{\vspace{5pt}{\large\color{deep}\bfseries #1}\par\vspace{1pt}}
\newcommand{\idea}[1]{\textcolor{deep}{\textbf{Ключевая идея.}} #1}
\newcommand{\point}[1]{\textcolor{deep}{\bfseries #1}}
\pgfplotsset{lessonaxis/.style={
    width=11.8cm,height=9.3cm,
    axis lines=middle,
    axis line style={black!75,thick,-{Stealth[length=2mm]}},
    tick align=outside,
    xtick distance=1,ytick distance=1,
    tick label style={font=\scriptsize,fill=white,inner sep=1.1pt},
    label style={font=\small},
    grid=major,grid style={black!13,thin},
    xlabel={$x$},ylabel={$y$},
    xlabel style={at={(axis description cs:1.03,0.52)},anchor=west},
    ylabel style={at={(axis description cs:0.5,1.03)},anchor=south},
    enlargelimits=false,clip=true,
    every axis plot/.append style={no markers},
}}
\newcommand{\answerline}{\par\vspace{7mm}\textcolor{textgray}{\footnotesize Место для решения}\par\vspace{4mm}\noindent\rule{\linewidth}{0.2pt}\par\vspace{8mm}\noindent\rule{\linewidth}{0.2pt}\par\vspace{8mm}\noindent\rule{\linewidth}{0.2pt}}
\begin{document}
\thispagestyle{empty}
\begin{titlepage}
\begin{center}
\vspace*{22mm}
{\large\color{deep} МАТЕМАТИКА}\par
\vspace{21mm}
{\LARGE\color{textgray} Чек-лист}\par
\vspace{14mm}
{\LARGE\color{accent} Квадратичная функция\\[4pt] и пересечение графиков}\par
\vspace{19mm}
{\large Определение коэффициентов по графику.\\ Координаты вершины. Узловые точки.}\par
\vfill
{\large Ксения Клыкова}\par
\vspace{11mm}
{\large 4 ноября 2026 г.}\par
\vspace{17mm}
{\normalsize Лёвин Артём Александрович\\[2mm]
 эксклюзивно для Ксении Клыковой}\par
\vspace{19mm}
\begin{tikzpicture}[x=1cm,y=1cm]
\draw[accent,line width=0.75pt] (0,0) .. controls (1.2,0.52) and (2.25,-0.50) .. (3.45,-0.07)
 .. controls (4.45,0.38) and (5.45,-0.32) .. (6.8,0.11);
\fill[accent] (3.45,-0.07) circle (1.1pt);
\end{tikzpicture}
\end{center}
\end{titlepage}
\setcounter{page}{1}

\lessonsec{1. Как узнать квадратичную функцию}
Квадратичная функция имеет вид
\[
 g(x)=ax^2+bx+c,\qquad a\ne0.
\]
Её график --- \textbf{парабола}. Каждый коэффициент отвечает за определённое свойство графика.

\begin{tabularx}{\linewidth}{@{}p{15mm}X@{}}
\toprule
\textbf{Что} & \textbf{Что можно узнать по графику}\\
\midrule
$c$ & Ордината пересечения с осью $Oy$: $c=g(0)$.\\
$a$ & Направление ветвей: $a>0$ --- вверх, $a<0$ --- вниз. При известной вершине величину $a$ можно определить сдвигом на 1 по $x$.\\
$b$ & Связан с абсциссой вершины: $x_{\text{в}}=-\dfrac{b}{2a}$.\\
\bottomrule
\end{tabularx}

\lessonsec{2. Вершина параболы}
Если $h=x_{\text{в}}$ и $k=y_{\text{в}}$, то
\[
 h=-\frac{b}{2a},\qquad b=-2ah,
 \qquad g(x)=a(x-h)^2+k.
\]
\idea{Вершина задаёт ось симметрии $x=h$. Точки на одинаковом расстоянии от этой прямой имеют одинаковые ординаты.}

\lessonsec{3. Приём «одна клетка вправо»}
Для формы $g(x)=a(x-h)^2+k$ получаем:
\[
 g(h+1)=a+k,\qquad\boxed{\ a=g(h+1)-k\ }.
\]
От вершины перейдите на \textbf{одну единицу вправо} по оси $x$. Изменение ординаты до параболы и есть коэффициент $a$ со знаком. При ветвях вниз изменение отрицательно.

\textbf{Ограничение практического приёма.} Формула верна и при дробном $h$. Но если вершина или нужная точка расположена между линиями сетки, её точные координаты могут быть нечитаемы. Тогда надёжнее воспользоваться узловыми точками.

\lessonsec{4. Пример: найдите формулу по вершине и точке}
Вершина $V(1;-4)$, парабола проходит через $(0;-3)$. Поэтому
\[
 g(x)=a(x-1)^2-4,\qquad -3=a-4\quad\Longrightarrow\quad a=1.
\]
\textbf{Итог:} $\boxed{g(x)=x^2-2x-3}$. Проверка: $g(0)=-3$.

\clearpage
\lessonsec{5. Метод узловых точек}
Когда вершина расположена неудобно, считайте координаты \textbf{трёх точек} графика и подставьте их в общую формулу:
\[
\left\{\begin{aligned}
 y_1&=ax_1^2+bx_1+c,\\
 y_2&=ax_2^2+bx_2+c,\\
 y_3&=ax_3^2+bx_3+c.
\end{aligned}\right.
\]
Получится система трёх линейных уравнений относительно $a,b,c$. Если есть точка при $x=0$, коэффициент $c$ известен сразу, и останется решить два уравнения.

\textbf{Пример.} Парабола проходит через $(-2;1)$, $(0;-3)$ и $(1;1)$.
Из $g(0)=-3$ получаем $c=-3$. Для двух других точек:
\[
 \begin{cases}4a-2b-3=1,\\a+b-3=1\end{cases}
 \Longleftrightarrow
 \begin{cases}2a-b=2,\\a+b=4.\end{cases}
\]
Сложив уравнения, получим $3a=6$, поэтому $a=2$, $b=2$.
\[
\boxed{g(x)=2x^2+2x-3.}
\]
Здесь $h=-\frac12$; точные координаты вершины необязательно считывать с рисунка.

\lessonsec{6. Как найти точку пересечения}
Если на одном рисунке изображены $y=f(x)$ и $y=g(x)$, то общие точки удовлетворяют равенству
\[
\boxed{f(x)=g(x).}
\]
\begin{enumerate}
\item Восстановите неизвестную формулу, если она задана только графиком.
\item Приравняйте выражения и найдите все значения $x$.
\item \textbf{Для нахождения ординаты} подставьте нужный $x$ в $f(x)$ или $g(x)$.
\end{enumerate}
\textbf{Пример.} Для $f(x)=5x+9$ и $g(x)=x^2+x-3$:
\[
5x+9=x^2+x-3\ \Longleftrightarrow\ x^2-4x-12=0
\ \Longleftrightarrow\ (x-6)(x+2)=0.
\]
Абсциссы пересечений $x=-2$ и $x=6$. Ординаты:
\[
f(-2)=-1,\qquad f(6)=39.
\]
\textbf{Ответ:} $(-2;-1)$ и $(6;39)$.

\lessonsec{7. Как различить две параболы}
Если одна формула известна, сначала найдите её $c$: график должен проходить через $(0;c)$. Дополнительно проверьте знак $a$ (направление ветвей) и положение вершины. Если обе параболы проходят через $(0;c)$, понадобятся дополнительные точки.

\textbf{Иллюстрация.} Известна формула $f(x)=x^2-2x+2$. Ей соответствует \textit{сплошной} график: он пересекает $Oy$ в точке $(0;2)$. У пунктирной параболы $c=-1$.
\begin{center}
\begin{tikzpicture}
\begin{axis}[lessonaxis,width=10.8cm,height=7.4cm,xmin=-1,xmax=3,ymin=-3.5,ymax=5.2]
 \addplot[accent,very thick,domain=-1:3,samples=150]{x^2-2*x+2};
 \addplot[second,very thick,dashed,domain=-0.8:2.8,samples=150]{-x^2+2*x-1};
 \addplot[only marks,mark=*,mark size=1.7pt,black] coordinates {(0,2) (0,-1)};
\end{axis}
\end{tikzpicture}
\end{center}

\lessonsec{8. Типичные ошибки}
\begin{itemize}[leftmargin=6mm,itemsep=2pt,topsep=1pt]
\item Путать $x$ в записи $g(x)$ и фиксированное число $h=x_{\text{в}}$.
\item Считать $b=-2h$: обязательно учитывать множитель $a$.
\item Принимать приблизительно считанные с рисунка координаты за точные.
\item Восстанавливать уже заданную в условии формулу повторно.
\item Завершать задачу на поиск \emph{ординаты} после нахождения только $x$.
\end{itemize}

\clearpage
\thispagestyle{plain}
\begin{center}
{\Large\color{deep}\bfseries Алгоритм восстановления формулы и пересечения}\par
\vspace{15mm}
\begin{tikzpicture}[
  node distance=13mm and 20mm,
  >={Stealth[length=2.7mm]},
  start/.style={ellipse,draw=deep,very thick,align=center,minimum width=45mm,minimum height=10mm,font=\small},
  phase/.style={rounded corners=3pt,draw=accent,thick,align=center,text width=92mm,minimum height=15mm,inner sep=5pt,font=\small},
  decision/.style={diamond,aspect=2.55,draw=deep,thick,align=center,text width=55mm,inner sep=2pt,font=\small},
  arr/.style={->,very thick,draw=deep}
]
\node[start] (s) {Изучите график};
\node[phase,below=of s] (c) {Определите $c=g(0)$; посмотрите на вершину и направление ветвей};
\node[decision,below=of c] (d) {Можно точно считать вершину и соседнюю точку?};
\node[phase,below left=17mm and 11mm of d,text width=43mm] (yes) {Найдите $a$ по сдвигу на 1; затем $b=-2ah$};
\node[phase,below right=17mm and 11mm of d,text width=43mm] (no) {Возьмите узловые точки и решите систему на $a,b,c$};
\node[phase,below=46mm of d] (formula) {Составьте $g(x)=ax^2+bx+c$ и проверьте по точкам};
\node[phase,below=of formula] (eq) {При поиске пересечений решите $f(x)=g(x)$};
\node[start,below=of eq] (finish) {Найдите нужную ординату $y$};
\draw[arr] (s)--(c);
\draw[arr] (c)--(d);
\draw[arr] (d.west) -| node[pos=.25,above,font=\footnotesize]{да} (yes.north);
\draw[arr] (d.east) -| node[pos=.25,above,font=\footnotesize]{нет} (no.north);
\draw[arr] (yes.south) |- (formula.west);
\draw[arr] (no.south) |- (formula.east);
\draw[arr] (formula)--(eq);
\draw[arr] (eq)--(finish);
\end{tikzpicture}
\end{center}

\clearpage
\section*{\color{deep}Тренировочный блок --- Путь А}
\textcolor{textgray}{Пять заданий с координатными графиками. Во всех рисунках одна клетка по осям соответствует одной единице. Считывайте только узловые точки, координаты которых определяются однозначно.}
\vspace{4mm}
\lessonsec{Задание 1. Вершина и свободный коэффициент}
На рисунке изображена парабола $g(x)=ax^2+bx+c$.
Найдите \textbf{все три коэффициента} $a,b,c$, запишите формулу функции и определите нули функции.
\begin{center}
\begin{tikzpicture}
\begin{axis}[lessonaxis,xmin=-2.2,xmax=4.2,ymin=-5.2,ymax=5.2]
 \addplot[accent,very thick,domain=-2:4,samples=240]{x^2-2*x-3};
 \addplot[only marks,mark=*,mark size=1.75pt,deep] coordinates {(1,-4) (0,-3)};
 \node[anchor=south west,font=\scriptsize,fill=white,inner sep=1pt] at (axis cs:1,-4) {$V$};
\end{axis}
\end{tikzpicture}
\end{center}
\answerline

\clearpage
\lessonsec{Задание 2. Ветви направлены вниз}
На рисунке изображена парабола $g(x)=ax^2+bx+c$.
Восстановите её формулу и вычислите $g(1)$.
\begin{center}
\begin{tikzpicture}
\begin{axis}[lessonaxis,xmin=-3.3,xmax=1.3,ymin=-6.2,ymax=4.2]
 \addplot[accent,very thick,domain=-3.12:1.12,samples=240]{-2*x^2-4*x+1};
 \addplot[only marks,mark=*,mark size=1.75pt,deep] coordinates {(-1,3) (0,1)};
 \node[anchor=south west,font=\scriptsize,fill=white,inner sep=1pt] at (axis cs:-1,3) {$V$};
\end{axis}
\end{tikzpicture}
\end{center}
\answerline

\clearpage
\lessonsec{Задание 3. Пересечение прямой и параболы}
На рисунке изображены прямая $f(x)=x+1$ \textcolor{second}{(пунктир)} и парабола $g(x)=ax^2+bx+c$ \textcolor{accent}{(сплошная)}. Восстановите формулу $g(x)$, найдите \textbf{координаты обеих} точек пересечения. Определите ординату точки $B$ с большей абсциссой.
\begin{center}
\begin{tikzpicture}
\begin{axis}[lessonaxis,xmin=-1.4,xmax=3.55,ymin=-1.3,ymax=6.3]
 \addplot[accent,very thick,domain=-1.25:3.45,samples=240]{x^2-2*x+1};
 \addplot[second,thick,dashed,domain=-1.4:3.55,samples=100]{x+1};
 \addplot[only marks,mark=*,mark size=1.85pt,black] coordinates {(0,1) (3,4)};
 \node[anchor=north west,font=\small,fill=white,inner sep=1pt] at (axis cs:0,1) {$A$};
 \node[anchor=south east,font=\small,fill=white,inner sep=1pt] at (axis cs:3,4) {$B$};
\end{axis}
\end{tikzpicture}
\end{center}
\answerline

\clearpage
\lessonsec{Задание 4. Вершина между узлами сетки}
Дана парабола $g(x)=ax^2+bx+c$ с вершиной, расположенной между вертикальными линиями сетки. Используя \textbf{три узловые точки}, найдите $a,b,c$, запишите формулу и решите уравнение $g(x)=1$.
\begin{center}
\begin{tikzpicture}
\begin{axis}[lessonaxis,xmin=-2.6,xmax=1.6,ymin=-4.3,ymax=4.3]
 \addplot[accent,very thick,domain=-2.35:1.35,samples=240]{2*x^2+2*x-3};
 \addplot[only marks,mark=*,mark size=1.85pt,deep] coordinates {(-2,1) (0,-3) (1,1)};
\end{axis}
\end{tikzpicture}
\end{center}
\answerline

\clearpage
\lessonsec{Задание 5. Две параболы}
На рисунке даны графики двух функций. Известно, что одна из них
\[
f(x)=2x^2-8x+5,
\]
а другая имеет вид $g(x)=ax^2+bx+c$. Определите, \textbf{какой график соответствует $f$} (сплошной или пунктирный). Восстановите $g(x)$ по узловым точкам и найдите \textbf{ординату точки $B$} --- правой точки пересечения графиков.
\begin{center}
\begin{tikzpicture}
\begin{axis}[lessonaxis,xmin=-0.5,xmax=4.55,ymin=-4.15,ymax=7.25]
 \addplot[accent,very thick,domain=-0.22:4.24,samples=240]{2*x^2-8*x+5};
 \addplot[second,very thick,dashed,domain=-0.5:4.55,samples=240]{x^2-3*x+1};
 \addplot[only marks,mark=*,mark size=1.85pt,black] coordinates {(1,-1) (4,5)};
 \addplot[only marks,mark=*,mark size=1.5pt,second] coordinates {(0,1) (2,-1)};
 \node[anchor=north east,font=\small,fill=white,inner sep=1pt] at (axis cs:1,-1) {$A$};
 \node[anchor=south east,font=\small,fill=white,inner sep=1pt] at (axis cs:4,5) {$B$};
\end{axis}
\end{tikzpicture}
\end{center}
\answerline

\clearpage
\section*{\color{deep}Ответы к тренировочному блоку}
\renewcommand{\arraystretch}{1.55}
\begin{tabularx}{\linewidth}{@{}p{13mm}X@{}}
\toprule
\textbf{№} & \textbf{Ответ}\\
\midrule
1 & $a=1$, $b=-2$, $c=-3$; $g(x)=x^2-2x-3$; нули: $x=-1$, $x=3$.\\
2 & $a=-2$, $b=-4$, $c=1$; $g(x)=-2x^2-4x+1$; $g(1)=-5$.\\
3 & $g(x)=x^2-2x+1$; $A(0;1)$, $B(3;4)$; ордината $B$ равна $4$.\\
4 & $a=2$, $b=2$, $c=-3$; $g(x)=2x^2+2x-3$; $x=-2$ или $x=1$.\\
5 & $f$ соответствует \textit{сплошной} линии; $g(x)=x^2-3x+1$; $B(4;5)$, ордината $B$ равна $5$.\\
\bottomrule
\end{tabularx}
\vspace{8mm}
\lessonsec{Итоговая памятка}
\begin{tabularx}{\linewidth}{@{}p{48mm}X@{}}
\toprule
\textbf{Что известно} & \textbf{Что делать}\\
\midrule
Пересечение с $Oy$ & Взять $c=g(0)$.\\
Точная вершина $(h;k)$ & Записать $g(x)=a(x-h)^2+k$.\\
Точка на 1 вправо от вершины & Вычислить $a=g(h+1)-k$.\\
Несколько узловых точек & Подставить их в $y=ax^2+bx+c$.\\
Две функции на рисунке & Решить $f(x)=g(x)$; затем найти нужное $y$.\\
\bottomrule
\end{tabularx}
\vfill
{\small\color{textgray}\centering Лёвин Артём Александрович эксклюзивно для Ксении Клыковой\par}
\end{document}
